\section{四、项目背景与政策}

\subsection{1.现实背景：黑土地退化危机与传统监测短板}

东北黑土地是我国重要的商品粮基地，承担着保障国家粮食安全的核心使命。黑土是历经漫长植被腐殖质积累而形成的珍贵土壤资源，其发育过程极为缓慢，一旦遭到破坏，在自然条件下短期内难以恢复。东北黑土区凭借肥沃的土壤条件与适宜的气候禀赋，长期支撑着我国大规模的粮食商品化生产，是国家粮食安全体系中不可替代的“压舱石”。

然而，长期高强度利用与不合理耕作，叠加自然侵蚀，导致黑土地出现严重退化问题，其表现可归纳为“变薄、变瘦、变硬”三个相互叠加的维度。就“变薄”而言，黑土层年均流失0.5-1厘米，部分区域耕作层不足20厘米，有效土层持续减薄，直接削弱了耕地的生产潜力（前述累计减薄6.83cm即为该速率多年累积的结果）；就“变瘦”而言，土壤有机质含量较开垦初期下降30\%-50\%，土壤的养分供给能力与结构稳定性明显降低；就“变硬”而言，长期翻耕与机械碾压造成土壤板结，土壤孔隙结构遭到破坏，保水保肥能力大幅下降。三者互为因果：土层变薄削弱了保肥基础，有机质下降又加剧了团粒结构解体，最终表现为耕地产能对投入的依赖不断增强。若不加以有效干预，将进一步影响区域粮食产能的可持续性。

三类退化的驱动机制存在本质差异，对应的治理路径与监测指标也不同。“变薄”以物理迁移过程为主导，表层土壤在水蚀、风蚀与耕作搬运的共同作用下持续流失，变化集中体现于地表裸露度、侵蚀沟发育与土层边界的空间位移；“变瘦”以生物化学过程为主导，有机质矿化速率长期超过腐殖质积累速率，体现为地表光谱特征与作物长势响应的整体性偏移；“变硬”则源于土壤结构退化，团粒结构解体、孔隙塌陷，导致降水入渗与持水能力下降。机理不同，治理手段便不能通用：以工程措施治理侵蚀难以解决有机质亏缺，以增施有机肥为主的培肥手段也无法修复已经塌陷的孔隙结构。这也说明，监测若只给出“是否变化”的二值结论，难以支撑分类施策。

与此同时，传统黑土地监测手段存在明显短板，难以支撑退化问题的及时发现与精准处置，具体体现在以下三个方面：

\begin{itemize}
    \item \textbf{人工巡查的覆盖瓶颈}：人工巡查方式依赖大量人力投入，成本高、周期长，且受制于交通条件与地块可达性，覆盖范围有限，难以实现大范围、高频次的连续观测，往往只能在退化显性化之后被动发现。
    \item \textbf{常规遥感的精度瓶颈}：常规遥感监测在亚米级精细尺度上的解译能力不足，受作物物候差异、光照条件、云雨天气与地表扰动等多重因素影响，解译结果中混入较多伪变化，误判率偏高，难以稳定区分真实的耕地退化与正常的农事活动。
    \item \textbf{多源数据的协同瓶颈}：多源数据分散在不同部门与业务系统中，标准不一、口径各异，存在明显的“数据孤岛”现象，难以形成“监测—预警—治理”的完整闭环，导致发现的问题无法及时转化为治理行动。
\end{itemize}

上述短板使得传统手段无法满足黑土地保护“早发现、早预警、早治理”的现实需求。因此，亟需构建一套天空地一体化、智能化、低成本的精准监测体系，以客观、连续的数据替代经验判断，实现退化风险的主动识别与前置干预。

本项目以国产高分二号亚米级光学卫星影像为主体数据源、以无人机影像为补充，正是针对上述瓶颈形成的组合方案：亚米级空间分辨率使地块尺度的地表覆盖差异得以分辨，为地块边界提取与地块内部变化定位提供几何基础；卫星平台的重访特性使“周级”观测频次成为可能，从而在时间维度上捕捉退化与恢复的渐变过程；无人机影像则弥补云雨天气与关键农时窗口的观测缺口。但真正决定监测价值的并非影像本身，而是能否从双时相影像中稳定识别真实的语义变化，这正是本项目以改进BIT模型为核心算法的原因。

\subsection{2.政策背景：国家顶层设计提供刚性支撑}

为应对黑土地退化危机，国家出台了一系列法律法规与重大政策，从法律约束、工程部署到配套措施形成了层层递进的制度链条，为本项目提供了明确的导向与坚实的依据。

从制度设计上看，三项政策各有侧重，共同构成“约束—目标—路径”的完整链条。《中华人民共和国黑土地保护法》解决的是“谁负责、保护什么”的问题，以法律形式确定保护范围与责任主体，约束力最为刚性；《国家黑土地保护工程实施方案（2021—2025年）》解决的是“达到什么水平、由谁投入”的问题，以量化目标与工程任务明确投入方向，兼具目标导向与资源配置功能；《东北黑土地保护性耕作行动计划（2020—2025年）》解决的是“用什么方式种地”的问题，以技术模式与作业规范引导基层生产行为。三者由刚性约束到柔性引导层层递进，其共同前提是保护成效必须可核查、可评价，否则责任追究、目标考核与落实核查都将缺乏事实依据。

\subsubsection{（1）核心法律保障}

2022年正式施行的《中华人民共和国黑土地保护法》，首次为特定土地类型制定专门法律，明确将黑龙江、吉林、辽宁、内蒙古四省区黑土耕地纳入法定保护对象，确立“因地制宜、用养结合、近期与远期目标结合”的基本原则，要求黑土地优先用于粮食、作物、蔬菜等农产品生产，除确无法修复的区域外，原则上全部纳入保护修复范围，为黑土地的系统性治理划定了法治红线。该法同时强化了地方政府的保护责任与监督检查要求，使黑土地保护由政策倡导上升为法定职责，也从制度层面为监测数据的规范采集、跨部门共享与业务化应用提供了依据。

\subsubsection{（2）重大工程支撑}

《国家黑土地保护工程实施方案（2021—2025年）》明确提出，到“十四五”末实现土壤有机质平均提高10\%以上，有效遏制黑土地退化趋势，重点支持天空地立体监测、精准监管、数字化治理技术研发与应用，与本项目的技术路线高度契合。方案将“用数据说话、靠技术管地”确立为重要实施方向，意味着监测能力建设不再是可有可无的辅助环节，而是工程成效可考核、可追溯的前提条件；只有把退化与恢复过程转化为可量化的空间信息，有机质提升等目标才具备客观的评价基础。

\subsubsection{（3）配套政策推动}

《东北黑土地保护性耕作行动计划（2020—2025年）》大力推广免耕少耕、秸秆还田等保护性耕作技术，国家重点研发计划也持续支持黑土地立体监测、智能诊断、精准管控技术的研发与落地，为项目实施提供了政策与资金层面的双重保障。保护性耕作会显著改变地表覆盖状态与耕作方式，其实际落实面积与实施效果难以通过传统上报方式准确掌握，同样依赖高频次、客观可信的遥感监测手段加以核查与校核，这与本项目所构建的“周级”动态监测能力形成了直接的衔接。

\begin{figure}[htbp]
    \centering
    \fbox{\parbox[c][4.5cm][c]{0.8\textwidth}{\centering 【图片空位：image3.png】东北黑土区保护性耕作面积逐年变化趋势示意图}}
    \caption{东北黑土区保护性耕作面积变化（2020—2025年）}
    \label{fig:2}
\end{figure}

\begin{figure}[htbp]
    \centering
    \fbox{\parbox[c][4.5cm][c]{0.8\textwidth}{\centering 【图片空位：image4.png】东北主要粮食作物调出价值变化趋势示意图}}
    \caption{东北主要粮食作物调出价值变化（2002—2022年）}
    \label{fig:3}
\end{figure}

\subsection{3.成效与需求：保护成效凸显，精准监测需求升级}

在一系列政策推动下，东北黑土地保护工作已取得阶段性成效。一方面，黑土区保护性耕作面积持续扩大，从2023年到2025年呈现稳步增长态势，免耕少耕等模式得到广泛推广，说明以少扰动、多覆盖为核心的耕作制度正在被基层生产主体接受并落地；另一方面，2002年至2022年间，东北主要粮食作物调出价值从不足500亿元增长至近3000亿元，印证了黑土地保护与粮食产能提升的协同效应，体现了黑土地作为国家粮食安全“压舱石”的关键作用。上述两方面变化相互印证，表明保护性耕作等举措在遏制退化趋势、稳定粮食产能方面已经产生了实际效果。

对上述两项趋势还需作进一步解读。保护性耕作面积的稳步扩大，反映的不只是推广力度的加强，更说明基层生产主体对少扰动、多覆盖耕作方式的主动接受程度在提高。粮食调出价值的持续增长则表明，黑土地在保护措施同步推进的条件下仍能稳定输出大规模商品粮，保护与产能之间并非此消彼长的对立关系。同时也应看到，面积数据主要来源于统计上报，价值变化还会受到市场价格波动等因素影响，二者均不能直接等同于耕地质量的实质改善；要判断有机质是否真正提升、黑土层是否真正增厚，仍须依靠直接观测获取的空间数据加以验证。

与此同时，保护工作的推进也使监测对象趋于复杂：地表覆盖形式更加多样，地块边界与耕作方式处于动态调整之中，退化与恢复过程往往在空间上零散分布、在时间上逐年渐变，仅凭人力与低精度手段难以准确刻画。当前，黑土地保护已进入“依法治土、精准治理”的新阶段，对监测技术的精准性、实时性、低成本性提出了更高要求：既要求能够分辨地块尺度的细微变化，又要求具备足够高的观测频次以捕捉变化过程，还必须把使用成本控制在基层可承受的范围内，否则技术能力难以真正下沉到县域与乡镇一线。

本项目聚焦黑土地保护的精准监测需求，正是为了进一步落实《黑土地保护法》及相关政策要求，通过智能化技术破解传统监测短板，助力实现黑土地“总量不减少、功能不退化、质量有提升、产能可持续”的保护目标，为国家粮食安全提供坚实的科技支撑。
